\section{Data Characterization: AEP/AJO Documentation and Workflows}
\label{sec:data}

This section empirically characterizes the corpora that underpin the causal-embedding
effort for Adobe Experience Platform (AEP) and Adobe Journey Optimizer (AJO). All
statistics below are computed directly from the repository's JSON/JSONL assets by the
reproducible script \texttt{implementation/analyze\_data.py}; the figures it emits are
included verbatim. We deliberately treat the data as small structured files (parsed with
\texttt{json}/\texttt{csv}) and compute everything with simple counting, so the numbers
are auditable and the analysis is decoupled from any model. We first describe the raw
sources (documentation pages and video tutorials), then the derived pairwise
causal-classification datasets and their tier system, then the labeling procedure
(adjacency and \texttt{MAX\_GAP} windowing) and the distribution shift it induces, and
finally the workflow-ordering and follow-up-reranking evaluation sets. We close with
implications for modeling.

\subsection{Source corpora}

\paragraph{Documentation.} The pairwise datasets are mined from Adobe Experience League
documentation. The deterministic phase-level extractor consumes a snapshot of $3{,}853$
scraped documentation pages (title/URL/content); after markdown, code-fence, table, and
callout cleaning, pages are split on H2/H3 headings into action-bearing
\emph{phase} paragraphs of $24$--$56$ words. Pages yielding between $3$ and $12$ phases
are retained, producing $1{,}301$ phase-structured workflows. A parallel sentence-level
builder restricts to \texttt{/tutorials/} pages and keeps only \emph{action} sentences
($5$--$60$ words containing an imperative verb such as \emph{select}, \emph{click},
\emph{navigate}, \emph{create}, \emph{configure}), dropping narration markers.

\paragraph{Video tutorials.} The transcript corpus
(\texttt{scrap\_transcript/aep\_transcripts.json}) contains $142$ AEP/AJO video tutorials,
each a \texttt{(title, url, toc\_path, video\_id, transcript)} record. Transcripts are
long and highly variable: the median length is $\approx 792$ words, the mean $\approx
1{,}135$ words, ranging from $175$ to $9{,}608$ words
(Figure~\ref{fig:transcripts}, left). All videos sit under the
\texttt{platform-learn/tutorials} tree; by topic area the largest buckets are
\emph{destinations} ($17$), \emph{ai-assistant} ($14$), \emph{schemas} ($13$),
\emph{privacy} ($11$), \emph{intro-to-platform} ($10$), and \emph{sources} ($10$)
(Figure~\ref{fig:transcripts}, right). The heavy right tail matters operationally: a
single transcript routinely exceeds a transformer's $512$-token budget, which is exactly
the ``text-starvation'' failure mode reported for cross-encoders on this data.

\begin{figure*}[t]
  \centering
  \includegraphics[width=\linewidth]{fig_transcript_lengths}
  \caption{Video-tutorial transcript corpus ($N=142$). Left: distribution of transcript
  length in words (median $\approx 792$, long right tail to $\sim 9.6$k words). Right:
  the twelve largest documentation topic areas. The corpus is AEP/AJO-centric and
  long-form, frequently exceeding standard $512$-token encoder budgets.}
  \label{fig:transcripts}
\end{figure*}

\subsection{Pairwise causal-classification datasets}

The central training signal is a \emph{directional pairwise} format: each row carries
\texttt{text\_1}, \texttt{text\_2}, and a binary \texttt{label}, where $\texttt{label}=1$
means \emph{\texttt{text\_1} causally precedes / enables \texttt{text\_2}}. The richer,
tiered variants additionally carry \texttt{tier\_1}, \texttt{tier\_2} (a curriculum-level
index), \texttt{sub\_1}, \texttt{sub\_2} (sub-chapter codes), and source URLs. We survey
the six datasets referenced as primary in the project context; their split sizes and
label balance are shown in Table~\ref{tab:datasets} and
Figures~\ref{fig:sizes}--\ref{fig:balance}.

\begin{table}[t]
  \centering
  \small
  \caption{Primary pairwise causal-classification datasets. Sizes are exact line counts;
  ``\%\,label\,1'' is the fraction of \emph{precedes} pairs in train; ``med.\ len.'' is the
  median word count of \texttt{text\_1}; ``tier'' marks datasets carrying tier metadata.
  Every dataset is essentially perfectly balanced by construction (forward/reverse pairs).}
  \label{tab:datasets}
  \begin{tabular}{lrrrcrc}
    \toprule
    Dataset & Train & Val & Test & \%\,lbl1 & med.\ len. & tier \\
    \midrule
    \texttt{aep\_causal\_cls34}    & 110{,}775 & 6{,}004  & 6{,}004  & 50.3 & 103 & \checkmark \\
    \texttt{hard\_neg\_sem}        & 261{,}210 & 14{,}736 & 6{,}004  & 50.3 & 103 & \checkmark \\
    \texttt{ajo\_doc\_not\_t1}     & 88{,}000  & 11{,}200 & 11{,}200 & 50.0 & 208 & \checkmark \\
    \texttt{new\_aep\_wf\_scrap}   & 13{,}904  & 1{,}676  & 1{,}548  & 50.0 & 44  & \\
    \texttt{workflowllm\_m2}       & 210{,}775 & 6{,}004  & 6{,}004  & 50.0 & 37  & \\
    \texttt{ajo\_newstyle}         & 20{,}897  & 2{,}612  & 2{,}613  & 50.0 & 95  & \\
    \bottomrule
  \end{tabular}
\end{table}

\paragraph{Label balance.} Because every dataset is built by emitting each ordered pair
in both directions ($\text{forward}\!\to\!1$, $\text{reverse}\!\to\!0$), the label
distribution is essentially exactly $50/50$ in every split
(Figure~\ref{fig:balance}). This is a deliberate and important property: it removes the
class-prior shortcut, forcing a model to learn \emph{directional asymmetry} rather than a
topical or frequency cue. It also means accuracy is directly interpretable against a
$50\%$ random floor, and that the genuine difficulty lives entirely in the
\emph{direction}, not in the marginal.

\paragraph{Scale and granularity.} Sizes span two orders of magnitude, from the
$\sim$17k-pair \texttt{new\_aep\_wf\_scrap} to the $\sim$282k-pair
\texttt{hard\_neg\_semantic} (Figure~\ref{fig:sizes}). Text granularity varies just as
widely: the workflow/phase datasets are short ($37$--$44$ word medians, single
action steps), whereas the conceptual document datasets are paragraph-scale
($103$--$208$ word medians). The \texttt{hard\_neg\_semantic} variant augments
\texttt{aep\_causal\_cls34} with semantically similar but causally unrelated distractors
(an explicit \texttt{hard\_neg} flag), roughly $2.4\times$ the base training size, to push
the model past lexical-overlap heuristics toward genuine causal discrimination.

\begin{figure}[t]
  \centering
  \includegraphics[width=\linewidth]{fig_dataset_sizes}
  \caption{Train/val/test sizes of the six primary pairwise datasets (stacked, exact line
  counts). Scale ranges from $\sim$17k to $\sim$282k pairs.}
  \label{fig:sizes}
\end{figure}

\begin{figure}[t]
  \centering
  \includegraphics[width=\linewidth]{fig_label_balance}
  \caption{Label balance (fraction of $\texttt{label}=1$ ``precedes'' pairs, train split).
  Forward/reverse pair construction yields a $\approx 50/50$ split everywhere, eliminating
  the class-prior shortcut.}
  \label{fig:balance}
\end{figure}

\subsection{The tier system and the cross-tier gap}

The tiered AEP datasets attach to each text a \emph{tier} index ($\mathrm{T}1$--$\mathrm{T}15$)
derived from a fixed sub-chapter$\to$tier mapping (\texttt{directional\_manifest.json}). The
tier is a coarse position in the AEP learning curriculum: low tiers are foundational
(schemas, ingestion), high tiers are downstream (activation, journeys). The tier gap
$\lvert \texttt{tier\_1} - \texttt{tier\_2} \rvert$ thus measures how conceptually far
apart the two texts are.

The gap distribution is the single most consequential property of these datasets
(Figure~\ref{fig:tiergap}). For \texttt{aep\_causal\_cls34} and
\texttt{hard\_neg\_semantic}, \emph{no} training pair has gap $0$: both texts always come
from \emph{different} tiers, and the gap is broadly spread (modes at $1$--$4$, a long tail
out to $14$), with adjacent-tier pairs (gap $1$) being only $\sim$$13.5\%$ of the data.
These datasets therefore teach \emph{cross-tier} precedence over a wide conceptual span.
In sharp contrast, \texttt{ajo\_doc\_not\_tier1} is built \emph{within} documents, so its
pairs are confined to gap $0$/$1$ (same or immediately adjacent tier). The two regimes are
complementary --- broad cross-tier supervision versus fine intra-document supervision ---
but they are \emph{not interchangeable}, and a model trained on one is evaluated
implicitly on the other.

\begin{figure*}[t]
  \centering
  \includegraphics[width=\linewidth]{fig_tier_gap_dist}
  \caption{Tier-gap $\lvert \texttt{tier\_1}-\texttt{tier\_2}\rvert$ distribution over
  training pairs. The cross-document datasets (\texttt{aep\_causal\_cls34},
  \texttt{hard\_neg\_semantic}) span gaps $1$--$14$ and contain \emph{no} same-tier pairs,
  whereas \texttt{ajo\_doc\_not\_tier1} is confined to gaps $0$/$1$. This complementarity
  is the structural source of train/test distribution shift.}
  \label{fig:tiergap}
\end{figure*}

\subsection{Adjacency / \texttt{MAX\_GAP} labeling and its distribution shift}

The workflow datasets reveal a second, finer adjacency constraint --- imposed not at the
tier level but at the \emph{step} level within a single document. The deterministic
extractor builds ordered pairs only between phases at most \texttt{MAX\_GAP} positions
apart: \texttt{MAX\_GAP}$=2$ (consecutive $+$ skip-one) for
\texttt{new\_aep\_workflow\_scrap}, and gap$\le 4$ (capped at $60$ pairs per document) for
the AJO sentence-level builder. Empirically, $100\%$ of \texttt{new\_aep\_workflow\_scrap}
training pairs have step gap $1$ ($56.5\%$) or $2$ ($43.5\%$) --- the model \emph{never}
sees a far-apart pair during training. Splits are document-level $80/10/10$ (seed $42$)
with no URL leakage across splits, and forward/reverse pairs keep labels balanced.

This windowing is well-intentioned --- it kills the topic shortcut and keeps pairs locally
coherent --- but it creates a hard \textbf{train/eval distribution shift}. At evaluation
time the orderer must score \emph{all} $\binom{n}{2}$ pairs of an $n$-step workflow,
including the far-apart (gap $\gg \texttt{MAX\_GAP}$) pairs that never appeared in
training. The independent-embedding bi-encoder, which excels on the adjacent in-distribution
pairs, degrades on these off-distribution long-range pairs; the cross-encoder's joint
cross-attention generalizes further, which is consistent with the repository's finding that
the DeBERTa-v3-large cross-encoder reaches $21/30$ correct workflow orderings on
\texttt{new\_ajo\_workflows} versus a bi-encoder's $0/30$. A second confound compounds it:
the format forces a \emph{total order} (every pair gets a direction) with \emph{no neutral
class}, yet real procedures are partial orders containing genuinely parallel/independent
steps. Forcing a chain over noisy, occasionally cyclic pairwise scores makes errors
compound. These two artifacts --- adjacency-only supervision and the total-order assumption
--- are the structural roots of the end-to-end ordering failure.

\subsection{Workflow-ordering and follow-up evaluation sets}

\paragraph{Workflow ordering.} The downstream ordering target
\texttt{ajo\_orchestrated\_workflows\_flat.json} contains $30$ orchestrated AJO
workflows, each \emph{uniformly} $3$ steps (\texttt{t1}/\texttt{t2}/\texttt{t3}), written
in a clean LLM-orchestrated style (e.g.\ select abandoned-cart profiles
$\to$ enrich with cart data $\to$ send personalized email). Four auxiliary
\texttt{aep\_3step\_workflows*} sets ($10$ workflows each; \emph{simple}, \emph{specific},
and \emph{transcript} stylistic variants) provide additional $3$-step probes. The small,
short, clean nature of these eval targets contrasts sharply with the long, noisy training
transcripts --- another axis of distribution shift between train and evaluation.

\paragraph{Follow-up reranking.} The follow-up dataset
(\texttt{prod\_jan\_feb\_26-aep-ajo\_follow-up-queries.json}) contains $334$ real AEP/AJO
anchor questions, each paired with \emph{exactly} $10$ candidate follow-up queries
(Figure~\ref{fig:followup}, left). Anchor questions are short --- mean $\approx 17.5$
words. The uniform $10$-candidate structure makes this a clean fixed-list reranking
benchmark (ideal for MRR / Recall@$k$), but it carries no gold relevance labels in the
raw file, so supervision or evaluation requires an external relevance signal.

\begin{figure}[t]
  \centering
  \includegraphics[width=\linewidth]{fig_followup_counts}
  \caption{Follow-up reranking set ($N=334$ anchor questions). Left: each question has
  exactly $10$ candidate follow-ups (a fixed-list reranking setup). Right: anchor-question
  length distribution (mean $\approx 17.5$ words).}
  \label{fig:followup}
\end{figure}

\subsection{Implications for modeling}

Four data properties should drive modeling decisions. \textbf{(1) Direction is the whole
task.} The enforced $50/50$ balance removes priors, so every gain must come from modeling
the antisymmetry $P(\text{A}\!\prec\!\text{B}) \neq P(\text{B}\!\prec\!\text{A})$ ---
motivating directional fusion, asymmetric scoring, or order-aware objectives over plain
symmetric similarity. \textbf{(2) Adjacency-only / \texttt{MAX\_GAP} supervision under-fits
long-range ordering.} Training only on gap-$\le$$2$/$4$ pairs but evaluating on all pairs
is the dominant distribution shift; closing it argues for training on the full transitive
closure within a document (far-apart pairs included), which also explains why the
cross-attention cross-encoder generalizes better than independent bi-encoder embeddings.
\textbf{(3) The total-order assumption is wrong.} With no neutral/abstain class, parallel
steps are mislabeled and forced into a chain; adding a neutral class and aggregating with a
cycle-tolerant ranker (Bradley--Terry / minimum-feedback-arc-set / Kemeny) rather than a
naive topological sort should reduce compounded ordering errors. \textbf{(4) Length and
style mismatch.} Long-form transcripts ($\sim$792-word median, $9.6$k-word tail) exceed
$512$-token budgets and differ stylistically from the short, clean $3$-step eval workflows;
this favors models that avoid truncating both texts into one shared budget and motivates
explicit train/eval style-gap mitigation. The tiered/hard-negative datasets and the fixed
$10$-candidate follow-up set are well-suited respectively to curriculum-style training and
to clean MRR-based reranking evaluation.
